\section{Method}
\label{sec:method}

\begin{figure*}[t]
\centering
\resizebox{0.82\textwidth}{!}{%
\begin{tikzpicture}[
  font=\small,
  box/.style={draw, rounded corners=2pt, minimum height=0.9cm, align=center, inner sep=4pt},
  fz/.style={box, fill=frozen!15, draw=frozen, thick},
  tr/.style={box, fill=trained!15, draw=trained, thick},
  tok/.style={draw, minimum width=0.38cm, minimum height=0.42cm, inner sep=0pt},
  arr/.style={-{Stealth[length=2.2mm]}, thick}]
  % inputs and the two streams
  \node[box, fill=gray!10, minimum width=1.9cm] (x) at (0,0) {Meeting\\audio $\bm{x}$};
  \node[fz, minimum width=2.7cm] (enc) at (3.2,1.0) {Audio encoder\\(Qwen3-Omni)};
  \node[fz, minimum width=2.7cm] (dia) at (3.2,-1.0) {Simultaneous diarizer\\(Nemotron 3), $S$ slots};
  \node[tr, minimum width=2.3cm] (proj) at (6.5,-1.0) {Speaker projector\\$\mathrm{Proj}_\psi$ (\S\ref{sec:method-thinker})};
  \draw[arr] (x.east) -- ++(0.3,0) |- (enc.west);
  \draw[arr] (x.east) -- ++(0.3,0) |- (dia.west);
  \draw[arr] (dia) -- node[above]{$\bm{P}$} (proj);
  % interleaved sequence
  \foreach \i/\c in {0/frozen,1/frozen,2/frozen,3/frozen,4/frozen,5/spkc,6/frozen,7/frozen,8/frozen,9/frozen,10/frozen,11/spkc} {
    \node[tok, fill=\c!35] (t\i) at (8.9+\i*0.4,0) {};
  }
  \node (dots) at (8.9+12*0.4+0.05,0) {$\cdots$};
  \draw[decorate, decoration={brace, amplitude=4pt}] (t0.north west) -- node[above=4pt, font=\small]{$\bm{A}_1$: $K$ acoustic tokens} (t4.north east);
  \node[below=2pt of t5, font=\small] {$\bm{s}_1$};
  \node[below=2pt of t11, font=\small] {$\bm{s}_2$};
  \node[font=\small] at (11.2,-1.25) {interleaved input $\bm{E}$ (Eq.~\ref{eq:interleave})};
  \draw[arr] (enc.east) -| node[pos=0.3, above, font=\small]{acoustic tokens $\bm{a}_i$} (t8.north);
  \draw[arr] (proj.east) -- node[above, font=\small]{$\bm{s}_j$} ++(0.75,0) |- (t0.west);
  % LLM
  \node[tr, minimum height=1.9cm, minimum width=2.0cm] (llm) at (15.6,0) {Thinker LLM\\+ LoRA\\(\S\ref{sec:method-thinker})};
  \draw[arr] (dots.east) -- (llm.west);
  % output segments (second row)
  \node[box, fill=white, anchor=west] (o1) at (6.2,-2.6)
    {$\spk{2}\,\tim{5}$\ \zh{我想知道会议的时间}\ $\tim{48}$};
  \node[box, fill=white, anchor=west] (o2) at (6.2,-3.7)
    {$\spk{3}\,\tim{42}$\ \zh{对，没错。}\ $\tim{56}$};
  \draw[arr] (llm.south) -- ++(0,-0.85) -| ([xshift=-0.3cm]o1.west) -- (o1.west);
  \draw[arr] ([xshift=-0.3cm]o1.west) |- (o2.west);
  \node[font=\small, anchor=west] at (6.2,-4.45) {who, when, what (Eq.~\ref{eq:serial})};
  % talker
  \node[tr, minimum width=2.4cm] (talk) at (14.9,-3.15) {PA Talker, voice $\bm{v}_k$\\one stream per speaker\\(\S\ref{sec:method-queue}--\ref{sec:method-train})};
  \draw[arr] (o1.east) -- ++(0.25,0) |- (talk.west);
  \draw[arr] (o2.east -| o1.east) ++(0,0) -- ++(0.25,0) |- (talk.west);
  \draw[thick] (o2.east) -- (o2.east -| o1.east);
  \node[box, fill=gray!10, minimum width=1.7cm] (y) at (17.6,-3.15) {Translated\\speech};
  \draw[arr] (talk) -- (y);
  % voice prompts
  \draw[arr, dashed, spkc!80!black] (dia.south) -- ++(0,-3.75) -| node[pos=0.25, below, font=\small]{voice prompt $\bm{v}_k$: solo regions of slot $k$ (\S\ref{sec:method-queue})} (talk.south);
  % simultaneous read loop
  \draw[arr, dotted, trained!80!black] (llm.north) -- ++(0,0.5) -- node[above, font=\small]{simultaneous: read every $\Delta$\,s, write or $\wait$ (\S\ref{sec:method-thinker})} ++(-7.4,0);
\end{tikzpicture}}
\caption{Overview of \sys{}. Acoustic tokens (what is said) are interleaved with speaker tokens that a trained projector computes from the slot activities $\bm{P}$ of a frozen simultaneous diarizer (who says it). Reading every $\Delta$ seconds, the Thinker writes speaker-first, time-stamped segments; the PA Talker speaks them in one stream per speaker and voice, trained on the posterior of the playback-alignment CRF, which does not run at test time. Blue: frozen; orange: trained.}
\label{fig:arch}
\end{figure*}


For every contribution to a meeting, \sys{} decides who said what (\S\ref{sec:method-thinker}) and when the listener
hears its translation (\S\ref{sec:method-queue}--\ref{sec:method-train}); Figure~\ref{fig:arch} gives an overview.

\subsection{Problem Formulation}
\label{sec:method-task}

Given a single-channel recording $\bm{x}\in\R^{L}$ of $T$ seconds ($L$ samples), \sys{} outputs segments
$\big(k_n,a_n,b_n,\bm{y}^{(n)}\big)_{n=1}^{N_\mathcal{Y}}$: a speaker label $k_n\in\{1,\dots,S\}$ that stays with one
person for the whole meeting, the start and end $a_n<b_n$ of the source speech in seconds, and the translation
$\bm{y}^{(n)}\in\mathcal{V}^{|\bm{y}^{(n)}|}$, a token sequence over the vocabulary $\mathcal{V}$. Segments are committed
at reads $t_1<\dots<t_R=T$ from $\bm{x}_{\le t_r}$ alone and never revised, and each is spoken in the output stream of
speaker $k_n$; the time at which a target unit (word or character) is played there is its \emph{heard time}, computed
from the playback schedule (\S\ref{sec:method-queue}), not measured on listeners.

\subsection{Speaker-Attributed Simultaneous Translation}
\label{sec:method-thinker}

Acoustic tokens carry no speaker identity across a meeting and are shared by overlapping speakers, so a speaker-blind
translator merges concurrent contributions (Figure~\ref{fig:motivation}). \sys{} therefore conditions the Thinker on a
speaker stream and has it name the speaker before translating.

\paragraph{Speaker tokens.} The frozen audio encoder of Qwen3-Omni~\citep{xu2025qwen3omni} maps $\bm{x}$ to acoustic
tokens $\bm{A}\in\R^{N\times d}$, $N$ tokens in the hidden space of size $d$ of its LLM, the Thinker. A frozen
simultaneous diarizer, Nemotron 3 Diarization~\citep{medennikov2026nemotron}, returns speaker-activity posteriors
$\bm{P}\in[0,1]^{F\times S}$ over $F$ frames and $S$ slots, which are ordered by first arrival and kept for the whole
meeting; overlapped speech activates several slots. For the $j$-th block of $K$ acoustic tokens, a trained projector
$\mathrm{Proj}_\psi$ maps the activities resampled to its token centres, $\bm{P}_j\in[0,1]^{K\times S}$, to a speaker
token $\bm{s}_j=\mathrm{Proj}_\psi(\bm{P}_j)\in\R^{d}$, placed after the block as in G-STAR~\citep{peng2026gstar}:
\begin{equation}
\bm{E}=[\bm{A}_1;\bm{s}_1;\dots;\bm{A}_M;\bm{s}_M]\in\R^{(N+M)\times d},
\label{eq:interleave}
\end{equation}
with $\bm{A}_j\in\R^{K\times d}$ and $M=\lceil N/K\rceil$; $\bm{E}$ fills the audio slot of the Thinker's prompt
(Appendix~\ref{app:train-thinker}).

\paragraph{Speaker-first, time-stamped output.} The output must then say who is speaking. We add a speaker token
$\spk{k}$ per slot and time tokens $\tim{m}$, marking $m\delta_t$ seconds from the window start, to $\mathcal{V}$, and
serialise a window's segments by start time, speaker first, as in serialised output training~\citep{kanda2020sot}:
\begin{equation}
\bm{y}=\mathop{\Vert}_{n=1}^{N_\mathcal{Y}}\big[\spk{k_n}\,\tim{\hat a_n}\,\bm{y}^{(n)}\,\tim{\hat b_n}\big]\in\mathcal{V}^{|\bm{y}|},
\label{eq:serial}
\end{equation}
where $\Vert$ concatenates and $\hat a_n,\hat b_n\in\mathbb{N}$ are $a_n,b_n$ in steps of $\delta_t$. The speaker token
conditions every translated word on its speaker, so two concurrent contributions become two segments; the time tokens
give each segment its source span. Targets use the diarizer's slots (each reference speaker is mapped to the slot most
active during their speech), so the Thinker learns to copy a slot rather than invent a label. The projector, the new
embeddings and LoRA adapters~\citep{hu2022lora} of the Thinker $p_\theta$ are trained with cross-entropy on $\bm{y}$,
speaker and time tokens up-weighted because one wrong label misattributes a whole sentence.

\paragraph{Reading and writing.} The Thinker reads every $\Delta$ seconds, $t_r=\min(r\Delta,T)$; the diarizer
finalises a frame $\gamma$ seconds after it, so read $r$ uses only frames up to $t_r-\gamma$. The Thinker continues its
committed output greedily until it emits $\wait$ ($\eos$ at the last read) and commits what precedes it. It learns when
to write from deadlines: each reference translation is split at punctuation into chunks, and chunk $l$, whose source
ends at $e_l$, is due at read $r(l)=\min\{r:t_r\ge e_l+\epsilon\}$ with interpreting lag $\epsilon$; the target of read
$r$ is the chunks due at $r$, serialised by Eq.~\eqref{eq:serial} and followed by $\wait$. This \emph{deadline Thinker}
is shared by \sys{} and its baseline B, so both commit the same text with the same speakers at the same times.

\subsection{Speaker Streams and Playback Queues}
\label{sec:method-queue}

Committed text is heard only after it is spoken, and one speaker's blocks are spoken in turn, so a stream that receives
more speech than its source time builds a queue; in B, two thirds of the delay after a source unit ends accrue after
its text is committed (Figure~\ref{fig:motivation}).

\paragraph{Voice-prompted streams.} The text a slot commits at one read forms a \emph{block}; one speaker's blocks of
the same read are merged (\emph{read-level merging}). The UniSS Talker~\citep{uniss2025} speaks each block as discrete
speech codes, conditioned on its text, the stream's earlier speech, a speaking-rate token and a voice prompt
$\bm{v}_k\in\R^{L_k}$, audio cut from regions where slot $k$ alone is active before its first commitment
(Appendix~\ref{app:train-talker}).

\paragraph{Queue and heard delay.} A speaker stream with blocks $i=1,\dots,I$ has release times $\bm{\nu}\in\R^{I}$ and
target-unit counts $\bm{W}\in\mathbb{N}^{I}$; a rendering $c$ of a block is a sequence of speech codes lasting $D(c)$
seconds. Block $i$ waits a \emph{backlog} $h_i\ge0$ (seconds) for the previous block to finish, with $h_1=0$ and
\begin{equation}
h_{i+1}=\max\big(0,\;\nu_i+h_i+D(c_i)-\nu_{i+1}\big).
\label{eq:backlog}
\end{equation}
Spreading its units evenly over its playback, the delay summed over block $i$'s units is the scalar
\begin{equation}
L_i(h_i,c_i)=W_i\,h_i+\tfrac{W_i+1}{2}\,D(c_i)
\label{eq:lat}
\end{equation}
in unit-seconds. A long rendering thus costs twice: its own units are spread over more time, and through
Eq.~\eqref{eq:backlog} every later block waits longer.

\paragraph{Backlog-aware speaking rate.} B's base Talker uses one fixed rate, so only a pause in the source drains a
queue. The playback layer starts from the block's
\emph{pace} $\tau^{\mathrm{pace}}_i$, the rate token at which its text would last as long as the source time since the
previous release, and raises it with the backlog:
\begin{equation}
\tau_i=\tau^{\mathrm{pace}}_i+\tau_0+\min\big(\bar\tau,\;\lfloor\kappa_h h_i\rceil\big),
\label{eq:rate}
\end{equation}
with offset $\tau_0$ and cap $\bar\tau$ in rate steps and $\kappa_h$ steps per second of backlog. An empty queue follows
the source's pace; a stream that falls behind speaks faster until its queue drains. The rule, however, reacts only to
the existing backlog: it cannot see that a long rendering chosen now delays every later block.

\subsection{Playback Alignment: A Latency-Constrained CRF}
\label{sec:method-plem}

Playback alignment (PA) charges each rendering also for the delay it imposes on later blocks by scoring whole playback
plans of a stream; since a rendering affects later blocks only through the backlog it leaves, a linear-chain
CRF~\citep{lafferty2001crf} over these plans admits exact inference.

\paragraph{Playback plans as a CRF.} For each block, the base Talker proposes a finite set $C_i$ of renderings, sampled
at the base rate and at a faster rate; the stream $\mathcal{U}=(\bm{\nu},\bm{W},C_{1:I})$ is the CRF's observation. A plan
$\bm{z}=(z_1,\dots,z_I)$ assigns block $i$ a state $z_i=(h_i,c_i)\in\mathcal{H}\times C_i$, where rounding
Eq.~\eqref{eq:backlog} up to a grid $\mathcal{H}=\{0,\delta_h,\dots,\bar H\}$ defines transitions
$g_i:\mathcal{H}\times C_i\to\mathcal{H}$, $h_{i+1}=g_i(h_i,c_i)$; $\bm{z}$ is feasible, $\bm{z}\in\mathcal{Z}(\mathcal{U})$,
if $h_1=0$ and it follows $g_i$. Each state receives the scalar score
\begin{equation}
\phi^\lambda_i(h,c)=\underbrace{\eta_\ell\,\bar\ell(c)+\eta_f\,f(c)}_{\text{quality}}-\underbrace{\lambda\,L_i(h,c)}_{\text{delay}},
\label{eq:plpot}
\end{equation}
where $\bar\ell(c)\le0$ is the base Talker's mean log-probability per code, $f(c)\le0$ the negative CTC loss per
character of the block's text under an independent recogniser~\citep{baevski2020wav2vec} (intelligibility),
$\eta_\ell,\eta_f>0$ fixed weights, and $\lambda\ge0$ the price of a unit-second. Normalising over feasible plans gives
\begin{equation}
p_\lambda(\bm{z}\mid\mathcal{U})=\frac{\mathds{1}[\bm{z}\in\mathcal{Z}(\mathcal{U})]}{Z_\lambda(\mathcal{U})}\exp\sum_{i=1}^{I}\phi^\lambda_i(z_i),
\label{eq:crf}
\end{equation}
with partition function $Z_\lambda(\mathcal{U})>0$: a linear-chain CRF whose node potentials are Eq.~\eqref{eq:plpot} and
whose pairwise potentials, $0$ or $-\infty$, encode the first-in-first-out queue. It is exact up to two approximations,
the sampled candidates and the backlog grid (Appendix~\ref{app:proofs}).

\paragraph{Exact inference.} Sums over the $\prod_i|C_i|$ candidate sequences can therefore be grouped by the backlog.
For candidate $c$ of block $i$, let $\bm{T}_i^{(c)}\in\{0,1\}^{|\mathcal{H}|\times|\mathcal{H}|}$ be its transition matrix,
$[\bm{T}_i^{(c)}]_{h,h'}=\mathds{1}[g_i(h,c)=h']$, and $\bm{\phi}_i^{(c)}\in\R^{|\mathcal{H}|}$ its potentials over $h$.
The forward and backward messages $\bm{\alpha}_i,\bm{\beta}_i\in\R_{\ge0}^{|\mathcal{H}|}$, whose entry $h$ is the total
weight of feasible prefixes leaving block $i$ with backlog $h$ and of all completions from there, satisfy
\begin{equation}
\begin{aligned}
\bm{\alpha}_{i+1}&=\textstyle\sum_{c\in C_i}\bm{T}_i^{(c)\top}\big(\bm{\alpha}_i\odot e^{\bm{\phi}_i^{(c)}}\big),\\
\bm{\beta}_i&=\textstyle\sum_{c\in C_i}e^{\bm{\phi}_i^{(c)}}\odot\bm{T}_i^{(c)}\bm{\beta}_{i+1},\\
\bm{\mu}_i^{(c)}&=\bm{\alpha}_i\odot e^{\bm{\phi}_i^{(c)}}\odot\bm{T}_i^{(c)}\bm{\beta}_{i+1}\,/\,Z_\lambda,
\end{aligned}
\label{eq:fb}
\end{equation}
with $\bm{\alpha}_1$ one-hot at $h=0$, $\bm{\beta}_{I+1}=\bm{1}$, $Z_\lambda=[\bm{\beta}_1]_0$, and $\odot$ the elementwise
product. $\bm{\mu}_i^{(c)}$ is column $c$ of the posterior marginal matrix $\bm{\mu}_i\in[0,1]^{|\mathcal{H}|\times|C_i|}$,
and the expected total delay $\Phi(\bm{z})=\sum_iL_i(z_i)$ of the stream is $\mathbb{E}_{p_\lambda}[\Phi]=\sum_{i,h,c}
\mu_i(h,c)L_i(h,c)$. As each $\bm{T}_i^{(c)}$ has one 1 per row, inference costs $O(I\,|\mathcal{H}|\max_i|C_i|)$ per
stream (Proposition~\ref{prop:fb}, Figure~\ref{fig:pa}).

\paragraph{A latency budget instead of a tuned price.} Tuning $\lambda$ by hand would need a sweep per direction, and its
scale has no direct meaning. We instead fix an interpretable budget, the expected delay per target unit, and let
$\lambda$ follow. Over a training corpus $\mathcal{S}$ of streams with $W_\mathcal{S}=\sum_{\mathcal{U}\in\mathcal{S}}\sum_i
W_i$ target units and summed delay $\Phi_\mathcal{S}$, we seek the distribution $q$ over the corpus's plans that stays
closest to the quality-only CRF $p_0$ within a budget of $b$ seconds per unit~\citep{ganchev2010posterior}:
\begin{equation}
\min_q\ \mathrm{KL}(q\,\Vert\,p_0)\quad\text{s.t.}\quad\mathbb{E}_q[\Phi_\mathcal{S}]\le bW_\mathcal{S},
\label{eq:pr}
\end{equation}
where $b$ and $W_\mathcal{S}$ are scalars. For a multiplier $\lambda\ge0$, the Lagrangian
$\mathrm{KL}(q\Vert p_0)+\lambda(\mathbb{E}_q[\Phi_\mathcal{S}]-bW_\mathcal{S})$ is minimised by the Gibbs distribution
$q_\lambda\propto p_0\,e^{-\lambda\Phi_\mathcal{S}}$. As $\Phi_\mathcal{S}$ is a sum of per-block delays, the factor
$e^{-\lambda\Phi_\mathcal{S}}$ only subtracts $\lambda L_i$ from every node potential, so
$q_\lambda=\prod_{\mathcal{U}\in\mathcal{S}}p_\lambda(\cdot\mid\mathcal{U})$: the price of delay in Eq.~\eqref{eq:plpot} is
the budget's Lagrange multiplier. Substituting $q_\lambda$ gives the dual function and its derivatives,
\begin{equation}
\begin{aligned}
\mathcal{D}(\lambda)&=-\textstyle\sum_{\mathcal{U}\in\mathcal{S}}\log\frac{Z_\lambda(\mathcal{U})}{Z_0(\mathcal{U})}-\lambda bW_\mathcal{S},\\
\mathcal{D}'&=\mathbb{E}_{p_\lambda}[\Phi_\mathcal{S}]-bW_\mathcal{S},\;
\mathcal{D}''=-\mathrm{Var}_{p_\lambda}\Phi_\mathcal{S}\le0,
\end{aligned}
\label{eq:dual}
\end{equation}
which yield the following guarantee.

\begin{proposition}[Latency budget]
\label{prop:pr}
If some distribution meets the budget of Eq.~\eqref{eq:pr} strictly, its solution is
$\prod_{\mathcal{U}\in\mathcal{S}}p_{\lambda^\star}(\cdot\mid\mathcal{U})$, where $\lambda^\star=0$ if
$\mathbb{E}_{p_0}[\Phi_\mathcal{S}]\le bW_\mathcal{S}$ and otherwise
$\mathbb{E}_{p_{\lambda^\star}}[\Phi_\mathcal{S}]=bW_\mathcal{S}$; the expected delay is non-increasing in $\lambda$.
\end{proposition}

\noindent The proposition makes the budget practical: the constrained posterior is the same CRF, computed by the same
pass of Eq.~\eqref{eq:fb}, and $\lambda^\star$ is found by bisection (proof and details in Appendix~\ref{app:proofs}).
We set $b$ below the delay of base-rate playback, so the budget binds.

\subsection{Training the Talker on the CRF Posterior}
\label{sec:method-train}

Because the CRF needs future release times, it runs only in training, where its posterior becomes the Talker's
training target.

\paragraph{Candidate weights.} On a training stream, we condition the marginals on the backlog $\tilde h_i\in\mathcal{H}$
that block $i$ meets under base-rate playback: row $\tilde h_i$ of $\bm{\mu}_i$, normalised, is the probability vector
$\bm{w}_i\in\Delta^{|C_i|}$ ($\Delta^n$ the probability simplex in $\R^n$) with entries
\begin{equation}
w_i(c)=e^{\phi_i(\tilde h_i,c)}\,\beta_{i+1}\big(g_i(\tilde h_i,c)\big)\big/\beta_i(\tilde h_i)
\label{eq:w}
\end{equation}
for $c\in C_i$, as $\alpha_i(\tilde h_i)$ cancels. Through $\beta_{i+1}$, a candidate is weighted by the quality and delay of every later
block it makes wait, so a rendering slightly better on its own but leaving a long backlog loses weight to a shorter one.

\paragraph{Objective.} The PA Talker $\pi_\xi$, initialised from the base Talker $\pi_0$, also reads a token
$\langle\tilde h_i\rangle$ of the binned backlog and is fine-tuned on the weighted candidates:
\begin{equation}
\begin{aligned}
\mathcal{J}(\xi)=\textstyle\sum_{\mathcal{U},i}\Big[&-\textstyle\sum_{c\in C_i}w_i(c)\log\pi_\xi(c\mid\cdot)\\
&+\kappa_\tau\,\mathrm{CE}\big(\bm{\rho}_i,\pi_\xi(\tau\mid\cdot)\big)\\
&+\kappa_{\mathrm{KL}}\,\mathrm{KL}\big(\pi_0\,\Vert\,\pi_\xi\big)\Big],
\end{aligned}
\label{eq:patrain}
\end{equation}
where $\pi_\xi(c\mid\cdot)$ is the probability of $c$'s codes given the block's text, voice prompt, history and backlog
token; $\bm{\rho}_i\in\Delta^{|\mathcal{T}|}$ is the posterior over the rate-token set $\mathcal{T}$,
$\rho_i(\tau)=\sum_cw_i(c)\mathds{1}[\tau(c)=\tau]$ with $\tau(c)$ the rate $c$ was sampled at, fitted by the
cross-entropy $\mathrm{CE}$; the KL keeps the codes close to $\pi_0$; and $\kappa_\tau,\kappa_{\mathrm{KL}}\ge0$ are weights.
Up to a constant, the first term equals $\mathrm{KL}(\bm{w}_i\Vert\tilde{\bm{\pi}}_\xi)-\log\pi_\xi(C_i)$, with
$\tilde{\bm{\pi}}_\xi\in\Delta^{|C_i|}$ the Talker renormalised on the candidates (Proposition~\ref{prop:patrain}): it fits
the Talker's preferences among renderings to the CRF posterior, analogous to one M-step of posterior-regularised
EM~\citep{ganchev2010posterior}. The target thus comes from a posterior that sees future release times, not from a
teacher network (cf.~\citealp{hu2016harnessing}).

\paragraph{Decoding.} At test time the PA Talker reads the backlog token of its actual queue, samples a few candidates
per block at the rate of Eq.~\eqref{eq:rate}, and keeps the one the recogniser finds most intelligible (highest $f$).
Read-level merging, the backlog-aware rate and the PA Talker form the \emph{playback layer}; \sys{} is the deadline
Thinker with this layer.

\subsection{cpASR-BLEU and tcpASR-BLEU}
\label{sec:method-metrics}

Inspired by cpWER~\citep{watanabe2020chime6} and tcpWER~\citep{neumann2023meeteval}, we propose cpASR-BLEU and tcpASR-BLEU,
which score the ASR transcripts of translated speech for content and speaker attribution, without and with time
constraints; we will release the scorer.

\paragraph{cpASR-BLEU.} For a meeting, let $\bm{o}_k$ be the transcript of output stream $k$ and $\bm{g}_\varrho$ the
reference translations of reference speaker $\varrho$, each concatenated in time order, and let $\sigma^\star$ be the
one-to-one assignment of streams to speakers that maximises the summed chrF~\citep{popovic2015chrf}, with $\mathcal{K}$
the matched streams. cpASR-BLEU is corpus BLEU~\citep{post2018call} over all meetings of the pairs
\begin{equation}
\begin{aligned}
\mathcal{P}_{\mathrm{cp}}={}&\{(\bm{o}_k,\bm{g}_{\sigma^\star(k)})\}_{k\in\mathcal{K}}\cup\{(\bm{o}_k,\varnothing)\}_{k\notin\mathcal{K}}\\
&\cup\{(\varnothing,\bm{g}_\varrho)\}_{\varrho\notin\sigma^\star(\mathcal{K})},
\end{aligned}
\label{eq:cpbleu}
\end{equation}
with $\varnothing$ the empty string, so a sentence spoken into the wrong stream costs twice: it is missing from its
speaker's pair and superfluous in another's.

\paragraph{tcpASR-BLEU.} tcpASR-BLEU adds the time constraint of tcpWER: every transcript segment, placed at its source
time with a collar, joins the overlapping unit $u$ (a speaker's consecutive utterances) of its matched speaker or
becomes an insertion. With $\bm{o}_u$ the segments joining $u$, $\bm{g}_u$ its reference and $\mathcal{I}$ the
insertions, it is corpus BLEU over $\mathcal{P}_{\mathrm{tcp}}=\{(\bm{o}_u,\bm{g}_u)\}_u\cup\{(\bm{o}_s,\varnothing)\}_{s\in\mathcal{I}}$,
so a correct translation given to the wrong speaker or time earns no credit (full definition: Appendix~\ref{app:eval}).
